\section{Positive Scalar Curvature}
\label{sec:psc}

\subsection{Overview}

brief definition, maybe intuition and osbtructions/backgroudn etc. based on overview paper

psc is somehow sweetspot for rigidity and flexibility

Probably the most important existence Theorem for positive scalar curvature metrics... Before stating the theorem let us briefly recall what surgery theory is about. The basic oberervation ist that a product $S^k \times S^{n-k-1}$ of spheres can be considered as the boundary of $S^k \times D^{n-k}$ or of $D^{k+1} \times S^{n-k-1}$... davor noch sowas schreiben wie "change the topology of a manifold in a controlled way"; remove thickened embedded $k$-sphere ($\to$ terminology of surgery of dimension k); EMBEDDINGS ALWAYS INTO INTERIOR (dann muss ich das unten in proposition zur surgery nicht mehr immer dazu schreiben) + can extend embedding of sphere iff trivial normal bundle; surgery intinmately connected to cobordism theory; how does surgery correspond to handles?; one common application of surgery theory is to "kill" certain homotopy groups $\to$ explain this process more precisely below 

\begin{theorem}[Schoen,Yau 1979, Gromov,Lawson,1980]
    \label{psc:thm:surgery_thm}
    (Surgery Theorem)
\end{theorem}

In order to apply this theorem we need a way to control the dimension of the surgeries performed or equivalently the indices of handles in a given cobordism. This question is completely independent of scalar curvature and has been subject of research especially in the context of the $h$-cobordims theorem. One central tool to control the indices of the handles in a cobordism is the following theorem due to Wall

\begin{theorem}[Wall, 1971]
    (Wall Theorem)
\end{theorem}

Combining gromov,lawson and Wall we see that: $M_0$ is psc, $M_1 \to W$ 2-connected then $M_1$ psc. In fact one can use the Theorem of Wall to obtain an even stronger result:

\begin{theorem}[extending psc to corbodism, Georg,Gaja]
    ...
\end{theorem}

So in view of the above discussion, the strategy to construct a psc metric on a given closed manifold $M_1$ of dimension at least $5$ is to first find a cobordism $W$ from some psc manifold $M_0$ to $M_1$ and then modify this cobordism by surgeries to obtain another cobordism $W^\prime$ from $M_0$ to $M_1$, such that the inlcusion $M_1 \to W^\prime$ is $2$-connected. Of course, such a modification by surgeries may not always be possible, so our next goal is to explain certain scenarios in which it is. In order to formulate these we recall the following definition. 

\begin{definition}[Spin and Totally non spin]
    ...
\end{definition}

We now have the following two existence theorems for psc metrics. One for the spin case and one for the totally non-spin case.

\begin{theorem}
    \label{psc:thm:existence_result_spin}
    Let $W^{n+1}$ be a cobordism from $M_0$ to $M_1$, $n \geq 5$, such that $W$ is spin. Suppose that $M_1$ is connected, the inlcusion $M_1 \to W$ induces an isomorphism on $\pi_1$ and that $M_0$ carries a psc-metric. Then $M_1$ also carries a psc-metric. In fact the psc-metric on $M_0$ can be extended to a psc-metric on $W$, which is of product form near the boundary. 
\end{theorem}

\begin{theorem}
    \label{psc:thm:existence_result_totally_non_spin}
    Let $W^{n+1}$ be an oriented cobordism from $M_0$ to $M_1$, $n \geq 5$. Suppose that $M_1$ is a connected totally non-spin manifold, the inlcusion $M_1 \to W$ induces an isomorphism on $\pi_1$ and that $M_0$ carries a psc-metric. Then $M_1$ also carries a psc-metric. In fact the psc-metric on $M_0$ can be extended to a psc-metric on $W$, which is of product form near the boundary. 
\end{theorem}

By the general strategy we laid out above, the proofs of these theorems will be using the given data to perform surgery on the cobordism $W$, such that the inlcusion $\iota \colon M_1 \to W$ becomes $2$-connected. By assumption $M_1$ is connected and $\iota$ already induces an isomorphism on $\pi_1$, so formulated differently our goal is to ``kill'' all generators in $\pi_2(W,M)$ by surgery. So as a first step let us make precise how one can use surgery to ``kill'' certain homotopy classes. 

\needspace{5cm}

\begin{proposition}
    \label{psc:prop:surger_killer}
    Let $W^n$ be smooth manifold with basepoint $w_0 \in W$, $M \subseteq \partial W$ a boundary component with basepoint $x_0 \in M$ and $2(k+1) \leq n$.
    \begin{enumerate}[label=(\roman*)]
        \item \label{psc:eq:surgery_killer_absolute_version} Let $\alpha_1,\ldots,\alpha_l \in \pi_k(W,w_0)$ be classes, which are represented by embeddings with trivial normal bundle. Then there exists a manifold $W^\prime$ obtained from $W$ by $l$-many $k$-surgeries away from the point $w_0$ together with an epimorphism
        \[
        \pi_k(W,w_0) \twoheadrightarrow \pi_k(W^\prime,w_0)
        \]
        containing the $\Z[\pi_1(W,w_0)]$-module generated by $\alpha_1,\ldots,\alpha_l$ in its kernel. Moreover 
        \[
        \pi_i(W,w_0) \cong \pi_i(W^\prime,w_0)
        \] 
        for all $i \leq k-1$.
        \item \label{psc:eq:surgery_killer_relative_version} Let $\alpha_1^\mathrm{rel},\ldots,\alpha_l^\mathrm{rel} \in \pi_k(W,M,x_0)$ be classes, which lift to classes $\alpha_1,\ldots,\alpha_l \in \pi_k(W,x_0)$ under the map $\pi_k(W,x_0) \to \pi_k(W,M,x_0)$ induced by the inclusion. Again assume that $\alpha_1,\ldots,\alpha_l$ are represented by embeddings with trivial normal bundle. Then there exists a manifold $W^\prime$ obtained from $W$ by $l$-many $k$-surgeries together with an epimorphism
        \[
        \pi_k(W,M,x_0) \twoheadrightarrow \pi_k(W^\prime,M,x_0)
        \]
        containing the $\Z[\pi_1(M,x_0)]$-module generated by $\alpha_1^\mathrm{rel},\ldots,\alpha_l^\mathrm{rel}$ in its kernel. Moreover 
        \[
        \pi_i(W,M,x_0) \cong \pi_i(W^\prime,M,x_0)
        \] 
        for all $i \leq k-1$.
    \end{enumerate}
\end{proposition}

\begin{proof}
    We first prove the absolute version \ref{psc:eq:surgery_killer_absolute_version} of the Proposition in three steps. \vspace{.5em}

    \noindent \textbf{Step 1.} We start by explaining how to kill a single homotopy class by surgery. Suppose that $w_0 \in W$ is an interior point. Let $\beta \in \pi_k(W,w_0)$ be a class represented by an embedding $f \colon (S^k,s_0) \to (W,w_0)$ for a specified basepoint $s_0 \in S^k$, which extends to an embedding $F \colon S^k \times D^{n-k} \to W$ into the interior of $W$. Choose any point $d_0 \in \partial D^{n-k} = S^{n-k-1}$ and a path $\lambda \colon [0,1] \to D^{n-k}$ from the center point $0 \in D^{n-k}$ to $d_0$. Let 
    \[
    w_0^\prime := F(s_0,d_0), \; \, \eta(t) := F(s_0,\lambda(t)).
    \]
    Then $\eta$ is a path from $w_0$ to $w_0^\prime$ and we can consider the change of basepoint isomorphism
    \[
    b_\eta \colon \pi_k(W,w_0) \to \pi_k(W,w_0^\prime)
    \]
    induced by $\eta$. We now perform the surgery. Let 
    \begin{align*}
        \tilde{W} &:= W - \mathrm{int}(\mathrm{im}(F))
        % W^\prime &:= \tilde{W} \cup_{S^k \times S^{n-k-1}} D^{k+1} \times S^{n-k-1}.
    \end{align*}
    and define $W^\prime$ by the pushout
    \[\begin{tikzcd}
	{S^k \times S^{n-k-1}} & {\tilde{W}} \\
	{D^{k+1} \times S^{n-k-1}} & {W^\prime}
	\arrow[from=1-1, to=1-2]
	\arrow[from=1-1, to=2-1]
	\arrow[from=1-2, to=2-2]
	\arrow[from=2-1, to=2-2]
    \end{tikzcd}\]
    as discussed above. Since the inclusion $S^k \times S^{n-k-1} \to D^{k+1} \times S^{n-k-1}$ is a cofibration it follows that the above suqare is a homotopy pushout, hence the map $\tilde{W} \to W^\prime$ has the same connectivity as $S^k \times S^{n-k-1} \to D^{k+1} \times S^{n-k-1}$, which is $k$-connected. Similiarly, the pushout 
    \[\begin{tikzcd}
	{S^k \times S^{n-k-1}} & {\tilde{W}} \\
	{S^k \times D^{n-k}} & W
	\arrow[from=1-1, to=1-2]
	\arrow[from=1-1, to=2-1]
	\arrow[from=1-2, to=2-2]
	\arrow[from=2-1, to=2-2]
    \end{tikzcd}\]
    implies that $\tilde{W} \to W$ is $n-k-1$-connected. In particular, the induced map $\pi_i(\tilde{W},w_0^\prime) \to \pi_i(W,w_0^\prime)$ is an isomorphism for $i \leq k$, since $2(k+1) \leq n$. Thus we can define maps
    \[
    \Phi_i \colon \pi_i(W,w_0) \xrightarrow{b_\eta} \pi_i(W,w_0^\prime) \xleftarrow{\cong} \pi_i(\tilde{W},w_0^\prime) \to \pi_i(W^\prime,w_0^\prime).
    \]
    By the above discussion $\Phi_i$ is an isomorphism for $i \leq k-1$ and an epimorphism for $i = k$. We now claim that $\Phi_k$ maps the class $\beta$ to $0$. First of all it is easy to see, that $b_\eta(\beta)$ may be represented by 
    \[
    \tilde{f} \colon (S^k,s_0) \to (W,w_0^\prime), \; y \mapsto F(y,d_0).
    \]
    Note that the image of $\tilde{f}$ lies in $\tilde{W}$ and thus $\Phi_i(\beta)$ is represented by
    \[
    \tilde{f} \colon (S^k,s_0) \to (W^\prime,w_0^\prime), \; y \mapsto \tilde{f}(y).
    \]
    This map factors as the following composition of inclusions
    \[
    S^k \to D^{k+1} \times \{d_0\} \to W^\prime,
    \]
    which is clearly contractible. Therefore $\Phi_k(\beta) = 0$. Note that all maps in the definition of $\Phi_k$ are compatible with the respective $\pi_1$-actions, hence the entire $\Z[\pi_1(W,w_0)]$-module generated by $\beta$ is contained in the kernel of $\Phi_k$. \vspace{.5em}

    \noindent \textbf{Step 2.} Next we explain how to kill multiple homotopy classes. To this end assume that $W$ is connected. Let $\beta_i \in \pi_k(W,w_i)$ for interior points $w_i \in W$, $1 \leq i \leq l$. Again suppose that $\beta_i$ can be represented by an embedding $f_i \colon (S^k,s_0) \to (W,w_i)$, which extends to an embedding $F_i \colon S^k \times D^{n-k} \to W$ into the interior of $W$. Moreover assume these embeddings to be disjoint, that is $\mathrm{im}(F_i) \cap \mathrm{im}(F_j) = \emptyset$ for $i \neq j$. This allows us to do surgery along all $F_i$ at once, i.e. we can define 
    \[
    \tilde{W} := W - \bigcup_{i=1,\ldots,l} \mathrm{int}(\mathrm{im}(F_i))
    \]
    and again let $W^\prime$ be defined by the pushout 
    \[\begin{tikzcd}
	{\coprod_{i=1,\ldots,l} S^k \times S^{n-k-1}} & {\tilde{W}} \\
	{\coprod_{i=1,\ldots,l}D^{k+1} \times S^{n-k-1}} & {W^\prime}
	\arrow[from=1-1, to=1-2]
	\arrow[from=1-1, to=2-1]
	\arrow[from=1-2, to=2-2]
	\arrow[from=2-1, to=2-2]
    \end{tikzcd}\]
    In order to make sense of the statement, that we can kill all homotopy classes $\beta_i$ at once we would like to consider them as elements in the same group. To achieve this let $d_0 \in S^{n-k-1}$ and $\lambda \colon [0,1] \to D^{n-k}$ be as in Step 1 and set 
    \[
    w_i^\prime := F_i(s_0,d_0), \; \, \eta_i(t) := F_i(s_0,\lambda(t)).
    \]
    Moreover choose any point $v_0 \in \tilde{W}$ together with paths $\gamma_i^\prime \colon [0,1] \to \tilde{W}$ from $w_i^\prime \to v_0$. Such paths exist, since $W$ was assumed to be connected and removing (thickened) spheres of codimension at least $2$ preserves connectedness. Let 
    \[
    \gamma_i := \gamma_i^\prime \ast \eta_i \colon [0,1] \to W
    \]
    denote the composed path from $w_i$ to $v_0$ and consider the following commutative diagram.
    \[\begin{tikzcd}
	{\pi_k(W,v_0)} & {\pi_k(W,v_0)} & {\pi_k(\tilde{W},v_0)} & {\pi_k(W^\prime,v_0)} \\
	{\pi_k(W,w_i)} & {\pi_k(W,w_i^\prime)} & {\pi_k(\tilde{W},w_i^\prime)} & {\pi_k(W^\prime,w_i^\prime)}
	\arrow["{\mathrm{id}}", from=1-1, to=1-2]
	\arrow["\cong"', from=1-3, to=1-2]
	\arrow[from=1-3, to=1-4]
	\arrow["{b_{\gamma_i}}", "\cong"', from=2-1, to=1-1]
	\arrow["{b_{\eta_i}}", "\cong"', from=2-1, to=2-2]
	\arrow["{b_{\gamma^\prime_i}}", "\cong"', from=2-2, to=1-2]
	\arrow["{b_{\gamma^\prime_i}}", "\cong"', from=2-3, to=1-3]
	\arrow["\cong"', from=2-3, to=2-2]
	\arrow[from=2-3, to=2-4]
	\arrow["{b_{\gamma^\prime_i}}", "\cong"', from=2-4, to=1-4]
    \end{tikzcd}\]
    As we have seen in Step 1, the class $\beta_i$ is contained in the kernel of the bottom horizontal composition, hence the same holds for the top horizontal compostion, which defines a map 
    \[
    \Phi \colon \pi_k(W,v_0) \to \pi_k(W^\prime,v_0).
    \]
    As in Step 1 we conclude that $\Phi$ is an epimorphism containing the $\Z[\pi_1(W,v_0)]$-module generated by $\beta_{\gamma_1}(\beta_1), \ldots, \beta_{\gamma_l}(\beta_l)$ in its kernel. Moreover, the same arguments as in Step 1 apply to show that $\pi_i(W,v_0) \cong \pi_i(W^\prime,v_0)$ for $i \leq k-1$. \vspace{.5em}

    \noindent \textbf{Step 3.} Finally we show how to kill multiple homotopy classes as formulated in \ref{psc:eq:surgery_killer_absolute_version}. Let $g_i \colon (S^k,s_0) \to (W,w_0)$ be embeddings with trivial normal  bundle representing the class $\alpha_i$. By a general position argument we can find embeddings $f_i \colon S^k \to W-\{w_0\}$ isotopic to $g_i$, such that $f_i(S^k) \cap f_j(S^k) = \emptyset$ for $i \neq j$. Moreover we may assume, that the image of $f_i$ is contained in the interior of $W$. Since the normal bundle of $g_i$ is trivial, the same holds for $f_i$. Therefore we can extend $f_i$ to an embedding 
    \[
    F_i \colon S^k \times D^{n-k} \to W - \{w_0\}
    \]
    into the interior of $W$, such that $\mathrm{im}(F_i) \cap \mathrm{im}(F_j) = \emptyset$ for $i \neq j$. Now we are in the setting as in Step 2 for $w_i = f_i(s_0)$ and $\beta_i = [f_i] \in \pi_k(W,w_i)$. Define paths $\gamma_i(t) := H_i(s_0,t)$. Then $b_{\gamma_i}(\beta_i) = \alpha_i$. We would like to apply Step 2 to $v_0 = w_0$, however the paths $\gamma_i$ are not of the same form as in Step 2 yet. We can however simply apply a homotopy to $\gamma_i$ to bring it into the desired form as follows. 

    ...

    Again let $W^\prime$ be the manifold obtained by surgery along $F_1 ,\ldots, F_l$ as before. Step 2 now yields an epimorphism
    \[
    \Phi \colon \pi_k(W,w_0) \to \pi_k(W^\prime,w_0).
    \]
    containing the $\Z[\pi_1(W,w_0)]$-module generated by the $\beta_{\tilde{\gamma}_i}(\beta_i) = \beta_{\gamma_i}(\beta_i) = \alpha_i$ in its kernel. Moreover $\pi_i(W,w_0) \cong \pi_i(W^\prime,w_0)$ for $i \leq k-1$. \vspace{.5em}

    \noindent \textbf{Step 4.} We can now deduce the relative version \ref{psc:eq:surgery_killer_relative_version} of the proposition using \ref{psc:eq:surgery_killer_absolute_version} applied to the given classes $\alpha_1, \ldots, \alpha_l \in \pi_k(W,x_0)$. Choose embeddings $F_i \colon S^k \times D^{n-k} \to W$ into the interior of $W$ as in Step 3 and let $\tilde{W}$ and $W^\prime$ be as before (see Step 2). For $i \leq k$ consider the following diagram, where all maps are induced by inlcusions and the rows are exact. 
    \[\begin{tikzcd}
	{\pi_i(M,x_0)} & {\pi_i(W,x_0)} & {\pi_i(W,M,x_0)} & {\pi_{i-1}(M,x_0)} & {\pi_{i-1}(W,x_0)} \\
	{\pi_i(M,x_0)} & {\pi_i(\tilde{W},x_0)} & {\pi_i(\tilde{W},M,x_0)} & {\pi_{i-1}(M,x_0)} & {\pi_{i-1}(\tilde{W},x_0)} \\
	{\pi_i(M,x_0)} & {\pi_i(W^\prime,x_0)} & {\pi_i(W^\prime,M,x_0)} & {\pi_{i-1}(M,x_0)} & {\pi_{i-1}(W^\prime,x_0)}
	\arrow[from=1-1, to=1-2]
	\arrow[from=1-2, to=1-3]
	\arrow[from=1-3, to=1-4]
	\arrow[from=1-4, to=1-5]
	\arrow["{\mathrm{id}}"', from=2-1, to=1-1]
	\arrow[from=2-1, to=2-2]
	\arrow["{\mathrm{id}}", from=2-1, to=3-1]
	\arrow["\cong"', from=2-2, to=1-2]
	\arrow[from=2-2, to=2-3]
	\arrow[from=2-2, to=3-2]
	\arrow[from=2-3, to=1-3]
	\arrow[from=2-3, to=2-4]
	\arrow[from=2-3, to=3-3]
	\arrow["{\mathrm{id}}"', from=2-4, to=1-4]
	\arrow[from=2-4, to=2-5]
	\arrow["{\mathrm{id}}", from=2-4, to=3-4]
	\arrow["\cong"', from=2-5, to=1-5]
	\arrow["\cong", from=2-5, to=3-5]
	\arrow[from=3-1, to=3-2]
	\arrow[from=3-2, to=3-3]
	\arrow[from=3-3, to=3-4]
	\arrow[from=3-4, to=3-5]
    \end{tikzcd}\]
    It follows from the Five Lemma that 
    \[
    \pi_i(\tilde{W},M,x_0) \to \pi_i(W,M,x_0)
    \] 
    is an isomorphism and 
    \[
    \pi_i(\tilde{W},M,x_0) \to \pi_i(W^\prime,M,x_0)
    \] 
    is an isomorphism for $i \leq k-1$. This shows $\pi_i(W,M,x_0) \cong \pi_i(W^\prime,M,x_0)$ for $i \leq k-1$. If $i = k$ the map $\pi_k(\tilde{W},x_0) \to \pi_k(W^\prime,x_0)$ is surjective and a diagram chase shows that 
    \[
    \pi_k(\tilde{W},M,x_0) \to \pi_k(W^\prime,M,x_0)
    \] 
    is surjective as well. Thus we can define the epimorphism
    \[
    \Phi_\mathrm{rel} \colon \pi_k(W,M,x_0) \xleftarrow{\cong} \pi_k(\tilde{W},M,x_0) \to \pi_k(W^\prime,M,x_0).
    \]
    We have already seen that $\alpha_1,\ldots,\alpha_l$ are contained in the kernel of the map 
    \[
    \pi_k(W,x_0) \xleftarrow{\cong} \pi_k(\tilde{W},x_0) \to \pi_k(W^\prime,x_0).
    \]
    So commutativity of the above diagram implies 
    \[
    \Phi_\mathrm{rel}(\alpha_i^\mathrm{rel}) = 0.
    \]
    Again all maps in the definition of $\Phi_\mathrm{rel}$ are compatible with the respective $\pi_1$-actions, such that $\Phi_\mathrm{rel}$ contains the entire $\Z[\pi_1(M,x_0)]$-module generated by $\alpha_1^\mathrm{rel},\ldots,\alpha_l^\mathrm{rel}$ in its kernel.

\end{proof}

Proof of proposition \ref{psc:prop:surger_killer} \\
Proof of Theorem \ref{psc:thm:existence_result_totally_non_spin} and \ref{psc:thm:existence_result_spin}

In the main construction in section ... we will invoke Theorem \ref{psc:thm:existence_result_totally_non_spin} to construct a psc metric, so let us conclude this section by discussing some basic results for identifying totally non-spin manifolds.

\begin{lemma}
    \label{lem:totally_non_spin_lemma} 
    \begin{enumerate}[label=(\roman*)]
        \item The product of a totally non-spin manifold with any non-empty manifold is totally non-spin. \label{eq:totally_non_spin_lemma1}
        \item A manifold containing a totally non-spin submanifold of codimension $0$ is totally non-spin. \label{eq:totally_non_spin_lemma2}
        \item If $M$ is totally non-spin and of dimension $n \geq 3$, then $M - D^n$ is totally non-spin for any embedded disk $D^n \subseteq \mathrm{int}(M)$ in the interior of $M$. \label{eq:totally_non_spin_lemma3}
    \end{enumerate}
\end{lemma}

\begin{remark}
    All statements of the above lemma remain valid if ``totally non-spin'' is replaced by ``non-spin''. The corresponding proofs are just simpler versions of the arguments given below.
\end{remark}

\begin{proof}%[Proof of Lemma \ref{lem:totally_non_spin_lemma}]
    For \ref{eq:totally_non_spin_lemma1} let $M_1$ and $M_2$ be manifolds with $M_1$ totally non-spin. Then the universal cover of the product $M_1 \times M_2$ is given by the product $\tilde{M_1} \times \tilde{M_2}$ of the respective universal covers. Denote by 
    \[
    \pi_1 \colon \tilde{M_1} \times \tilde{M_2} \to \tilde{M_1}, \; \;  \pi_2 \colon \tilde{M_1} \times \tilde{M_2} \to \tilde{M_2}
    \]
    the projections and let $\iota \colon \tilde{M_1} \to \tilde{M_1} \times \tilde{M_2}$ be an inclusion for some choice of basepoint in $\tilde{M_2}$. We compute 
    \begin{align*}
    \iota^\ast w_2(T(\tilde{M_1} \times \tilde{M_2})) & = \iota^\ast w_2(\pi_1^\ast T \tilde{M_1} \oplus \pi_2^\ast T \tilde{M_2}) \\
    & = w_2((\pi_1 \circ \iota)^\ast T\tilde{M_1} \oplus (\pi_2 \circ \iota)^\ast T\tilde{M_2} ) \\
    & = w_2(T \tilde{M_1}) \neq 0,
    \end{align*}
    where we used that $\pi_1 \circ \iota$ is the identity and that $\pi_2 \circ \iota$ is constant, such that $(\pi_2 \circ \iota)^\ast T\tilde{M_2}$ is trivial. Thus $w_2(T(\tilde{M_1} \times \tilde{M_2})) \neq 0$.

    For \ref{eq:totally_non_spin_lemma2} suppose that $N \subseteq M$ is some totally non-spin codimension $0$ submanifold of the manifold $M$. Let $\iota \colon N \to M$ be the inclusion and consider its lift to the universal covers:
    \[\begin{tikzcd}
	{\tilde{N}} & {\tilde{M}} \\
	N & M
	\arrow["{\tilde{\iota}}", from=1-1, to=1-2]
	\arrow[from=1-1, to=2-1]
	\arrow[from=1-2, to=2-2]
	\arrow["\iota", from=2-1, to=2-2]
    \end{tikzcd}\]
    Since $\tilde{\iota}$ is a local diffeomorphism, the differential $d \tilde{\iota} \colon T\tilde{N} \to T\tilde{M}$ is a fiber-wise isomorphism, such that the following is a pullback diagram
    \[\begin{tikzcd}
	{T\tilde{N}} & {T\tilde{M}} \\
	{\tilde{N}} & {\tilde{M}}
	\arrow["{d \tilde{\iota}}", from=1-1, to=1-2]
	\arrow[from=1-1, to=2-1]
	\arrow[from=1-2, to=2-2]
	\arrow["{\tilde{\iota}}", from=2-1, to=2-2]
    \end{tikzcd}\]
    Therefore $\tilde{\iota}^\ast w_2(T\tilde{M}) = w_2(T\tilde{N}) \neq 0$ and so $w_2(T\tilde{M}) \neq 0$.

    For \ref{eq:totally_non_spin_lemma3} assume without of loss generality that $M$ is connected. Denote by $B^n := \mathrm{int}(D^n)$ the open ball. Then $M - D^n$ is the interior of $M - B^n$ 
    from which we conclude that $M - D^n$ is totally non-spin if and only if $M- B^n$ is totally non-spin. This follows using an analogous argument as in \ref{eq:totally_non_spin_lemma2} by noting that for a manifold $X$ with boundary the universal cover of $\mathrm{int}(X)$ is given by $\mathrm{int}(\tilde{X})$ and the inlcusion $\mathrm{int}(\tilde{X}) \to \tilde{X}$ is a homotopy equivalence. We will now show that $M - B^n$ is totally non-spin. Let $p \colon \tilde{M} \to M$ be the universal cover of $M$ and write $M_0 := M - B^n$. We claim that the universal cover of $M_0$ agrees with the covering 
    \[
    \hat{M}_0 = \tilde{M}|_{M_0} = \tilde{M} - p^{-1}(B^n) \xrightarrow{p} M_0.
    \]
    %Für den Beweis, dass $\hat{M}_0$ connected ist, verwendet man im Grunde auch, dass $\hat{M}_0$ homotopieäquivalent zu $\tilde{M}-p^{-1}(0)$ ist, was wiederrum connected ist, da $n \geq 2$, sodass isolierte Punkte herausnehmen, nichts daran ändert, dass die Mannigfaltigkeit zusammenhängend ist. Also hier das ganze vllt noch leicht umstrukturieren.
    First of all $\hat{M}_0$ is connected as $n \geq 2$. We now have to show that $\hat{M}_0$ is simply connected. Let $F = p^{-1}(x)$ be some fiber of $p$. From the classification of coverings we see that 
    \[
    p^{-1}(B^n) \cong B^n \times F
    \]
    since $B^n$ is simply connected. If the fiber $F$ is finite we can apply the Seifert-van-Kampen Theorem to conclude that $\tilde{M}$ with a single disk $B^n$ removed is still simply connected, such that the statement follows inductively. If the fiber is not finite we can argue more generally as follows. Let $\phi \colon S^1 \to \hat{M}_0$ be some smooth loop. Certainly we can contract this loop in $\tilde{M}$, that is we can find some map $\Phi \colon D^2 \to \tilde{M}$ restricting to $\alpha$ on the boundary. We now show that $\Phi$ is homotopic rel$\,S^1$ to some map $\Phi_1 \colon D^2 \to \tilde{M}$, such that the image of $\Phi_1$ is already contained in $\hat{M}_0$, from which the claim follows. To construct such a homotopy let $0 \in B^n \subseteq M$ be the center of the embedded disk. After a preliminary homotopy we may assume that $\Phi$ is transversal to $p^{-1}(0)$. The crucial observation is that, since $n \geq 3$, this is equivalent to requiring that $\Phi$ does not intersect $p^{-1}(0)$. Now note that $S^{n-1}$ is a strong deformation retract of $D^n - 0$ and applying such a deformation retraction on each disk $D^n - 0 \subseteq p^{-1}(D^n - 0)$ shows that $\hat{M}_0$ is a strong deformation retract of $\tilde{M} - p^{-1}(0)$. Let 
    \[
    r \colon \tilde{M} - p^{-1}(0) \to \hat{M}_0
    \]
    be such a retraction and let 
    \[
    r_t \colon \tilde{M} - p^{-1}(0) \to \tilde{M} - p^{-1}(0)
    \]
    be a homotopy rel$\,S^1$ from the identity to $r$ (composed with the corresponding inclusion). Then $\Phi_t := r_t \circ \Phi$ defines a homotopy rel$\,S^1$ from $\Phi$ to some map $\Phi_1$ whose image is contained in $\hat{M}_0$. 
    % Since $\Phi_1^\prime(S^1) = \phi(S^1)$ is contained in the interior $\hat{M}_0$ of this manifold we can take a further homotopy rel$\,S^1$ from $\Phi_1^\prime$ to a map $\Phi_1$ whose image is contained in $\hat{M}_0$. 
    This shows the claim and hence the universal cover of $M_0$ is given by 
    \[
    \hat{M}_0 = \tilde{M} - p^{-1}(B^n) \cong \tilde{M} - (B^n \times F).
    \]
    Finally, the Mayer-Vietoris sequence of the pushout 
    \[\begin{tikzcd}
	{S^{n-1} \times F} & {D^n \times F} \\
	{\tilde{M} - (B^n \times F)} & {\tilde{M}}
	\arrow[from=1-1, to=1-2]
	\arrow[from=1-1, to=2-1]
	\arrow[from=1-2, to=2-2]
	\arrow["\iota", from=2-1, to=2-2]
    \end{tikzcd}\]
    where all maps are inclusions, then shows that $\iota$ is injective on $H^2(-;\Z_2)$. Therefore
    \[
    w_2(T\hat{M}_0) = w_2(\iota^\ast T\tilde{M}) = \iota^\ast w_2(T\tilde{M}) \neq 0.
    \]
\end{proof}

A standard example of a non-spin manifold is $\CP^2$. Since $\CP^2$ is simply connected this is also an example of a totally non-spin manifold, which we will appeal to in section ...

\begin{lemma} \label{lem:cp2_not_spin}
    $\C P^2$ is non-spin.
\end{lemma}

(proof)
